\originalchapter{声子、涨落与反常响应}\label{ch:phonon}
\originalsection{晶格正常模式与Debye热容}
晶格上的位移只在离散位置$\vct r$取值. 相差倒格矢$\vct G$的波矢满足$\e^{\ii(\vct k+\vct G)\cdot\vct r}=\e^{\ii\vct k\cdot\vct r}$, 因而给同一模式. 有$N_c$个原胞时第一Brillouin区含$N_c$个独立波矢, 每原胞有$s$个原子则每波矢有$3s$个振动分支. 对简单立方晶格$-\pi/a<k_i\le\pi/a$是一个可选区间; 不是所有晶体的Brillouin区都为立方体.

平移对称的二阶刚度矩阵只依赖$\vct r-\vct r'$. Fourier变换使不同$\vct k$解耦, 对每个波矢再对角化$3s\times3s$矩阵得到$\omega_j(\vct k)$. 一个分支的态密度满足$D(\omega)=V\int\delta(\omega-\omega(\vct k))\dd^3k/(2\pi)^3$; 色散曲线平坦处同一频带聚集更多态, 解释原课态密度图的峰. 数值数据的具体峰不是本稿用光滑示意所证明的结论.

Einstein模型把$3N$个频率都设为$\omega_E$, 热容为$3N$倍单振子热容, 低温指数下降. 实际固体有长波无隙声学模式$\omega_j\simeq c_jk$, 三维低频态密度
\[
D(\omega)=\frac{V\omega^2}{2\pi^2}(c_l^{-3}+2c_t^{-3}).
\]
将热激发能积分, 低温可把有限截止近似延至无穷, 被遗漏尾项指数小, 得
\[
U-U_0=\frac{\pi^2V\kb^4T^4}{30\hbar^3}(c_l^{-3}+2c_t^{-3}),\quad
C_V=\frac{2\pi^2V\kb^4T^3}{15\hbar^3}(c_l^{-3}+2c_t^{-3}).
\]
若用等效声速和截止$\omega_D$使总态数$3N$, 定义$\Theta_D=\hbar\omega_D/\kb$, 就得到$C_V\simeq12\pi^4N\kb(T/\Theta_D)^3/5$. 高温所有模经典化恢复$3N\kb$. 在$d$维线性色散中, 态数低能尺度为$k^d$, 所以$C\propto T^d$. 低温幂次来自无隙模式及其态密度, 并非每个模式的热容本身是幂函数.

\originalsection{零点能的压力与同位素效应}
声子基态能$U_{\rm zp}=\frac12\sum_{\vct k,j}\hbar\omega_j$虽不贡献定频热容, 却会随体积改变. 定义Gr\"uneisen参数$\gamma_j=-\partial\ln\omega_j/\partial\ln V$, 从$F=\phi(V)+\sum F_j$求$P=-F_V$, 得
\[
P=-\phi'(V)+\frac1V\sum_j\gamma_j\avg{E_j}.
\]
$\phi(V)$是静态晶格势能, 同频率尺度近似$\gamma_j=\gamma$时, 零外压平衡要求$\phi'(V)=\gamma U_{phonon}/V$. 所以即使$T=0$也会因零点能而偏离静态势的极小体积.

同位素主要改变原子质量, 相似力常数下$\omega\propto m^{-1/2}$. ${}^6$Li与${}^7$Li零点能之比约$\sqrt{7/6}$, 轻同位素零点膨胀更大. 原课给零温晶格间距差约0.07\%作为观察例; 这个百分数不能只从质量比算出, 还需要晶格刚度和$\gamma$. 在低温$U_{phonon}=U_{\rm zp}+aT^4$使热膨胀呈相应温度趋势.

\originalsection{热力学涨落与临界乳光}
正则能量涨落与热容相连, 巨正则粒子数涨落与压缩率相连. 用$\dd\mu=v\dd P$得$(\partial n/\partial\mu)_T=n^2\kappa_T$, 因而
\[
\Var N=\kb TVn^2\kappa_T,\qquad
\frac{\Var n}{n^2}=\frac{\kb T\kappa_T}{V}.
\]
有限压缩率下相对涨落随体积增大消失, 临界点附近$\kappa_T$变大时却可能显著. 密度起伏产生折射率起伏, 在许多尺度上散射光, 形成临界乳光. 小子体积的涨落可用这一体相关系, 但临界相关长度接近容器尺度时需处理有限尺寸, 不能直接令方差无穷.

\originalsection{熔化曲线与反常热膨胀}
稳定性要求$\kappa_T\ge0$, 因而均匀稳定相中加压不会增大体积. 但$\alpha$的符号没有同样约束, 只能由$C_P-C_V=TV\alpha^2/\kappa_T\ge0$约束它的平方. 水结冰的体积增大也不违反热力学; 它是两个不同稳定相的体积差, 不是某一相的负压缩率.

共存线上的$\mu_1=\mu_2$, 求微分得Clapeyron式
\[
\frac{\dd P}{\dd T}=\frac{s_2-s_1}{v_2-v_1}=\frac{L}{T(v_2-v_1)},
\]
$s,v$是每粒子熵和体积, $L$是由相1变到相2所吸热. 熵较大的相随升温更有利, 但“液体总比固体熵大”并非无条件. 低温${}^3$He固体中核自旋近自由, 可有每粒子约$\kb\ln2$的熵, 而Fermi液体低温熵正比$T$. 固体熵更大且体积更小的温区内熔化曲线斜率可以为负, 固定压强升温可发生反常凝固. 更低温自旋相互作用必须使固体熵趋第三律极限, 所以这个简化熵模型不许外推到严格零温. 后面的相图题将按原卷实际曲线分区判断, 不靠一般口号猜相.
